\subsection{Recovery of a prescribed volume limit}
\label{subsec:fixed-spacing-recovery}\label{sec:recovery}
An infinite-lattice law gives a step-field law by
\cref{prop:step-realization}. To treat a prescribed volume limit
of the finite-torus fields, we also need the converse
identification. We use the bump-function argument of
\cite[Section~6.4]{QED3Manuscript}.

Fix a spacing $a>0$. Let $\mathcal P_2$ be the positively oriented
plaquettes of $\Z^3$, and put
\[
 \Omega=\{u\in\R^{\mathcal P_2}:du=0\},
\]
with its coordinate product topology. For each integer $n\ge4$,
let $\lambda_{a,n}$ be a probability law of real closed
two-cochains on $(\Z/n\Z)^3$ satisfying
\begin{equation}\label{eq:recovery-finite-gaussian}
 \E_{\lambda_{a,n}}e^{\langle\beta,u\rangle}
 \le e^{\|\beta\|_{\ell^2((\Z/n\Z)^3)}^2/2}
 \qquad(\beta\text{ a real torus two-cochain}).
\end{equation}
Both the pairing and the norm sum over positively oriented torus
plaquettes. Denote the law of the periodic extension by
$\nu_{a,n}$. Use the centered interpolation and coefficients
$f_{[a]}$ of \eqref{eq:plaquette-cells}, and write $P_{a,n}$
for its law on $H_{-4}$.
For general Schwartz tests these finite-torus laws obey the
periodized bound \eqref{eq:limit-cutoff-G} below, with physical
period $R=an$. The following identification gives the sharp
$L^2$ bound after the volume limit.

\begin{proposition}\label{prop:fixed-spacing-recovery}
Suppose $n_j\to\infty$ and $P_{a,n_j}\Longrightarrow\mu_a$
weakly on $H_{-4}$. Then there is a probability law $\nu_a$
on $\Omega$ such that
\begin{equation}\label{eq:limit-volume-identification}
 \nu_{a,n_j}\Longrightarrow\nu_a,\qquad
 \mu_a=(\sharp_a)_*\nu_a,
\end{equation}
where $\sharp_a$ is the centered step map, defined as a Borel
map into $H_{-4}$ on a set of full $\nu_a$ measure and set
equal to zero elsewhere. The convergence of cochain laws is
in the product topology and holds along the given sequence.
For every finitely supported real cochain $\beta$,
\begin{equation}\label{eq:limit-lattice-G}
 \E_{\nu_a}e^{\langle\beta,u\rangle}
 \le e^{\|\beta\|_{\ell^2(\mathcal P_2)}^2/2}.
\end{equation}
In particular, $\mu_a$ satisfies the Gaussian bound
$\E_{\mu_a}e^{F(f)}\le e^{\|f\|_2^2/2}$ for real Schwartz
two-forms $f$.
\end{proposition}

\begin{proof}
Discard finitely many terms so that $an_j\ge1$.
The finite-volume step maps take values in $H_{-3}$.
Indeed, the periodization estimate below and the Hermite bound
$p_1(h_ke_\alpha)\le31\lambda_k$ from
\cite[Section~3]{QED3Integration} give
$|F_a(h_ke_\alpha)|\le31\lambda_k\|u\|_{\ell^2}$.
Summing with weight $\lambda_k^{-6}$ proves continuity of
each finite-dimensional step map into $H_{-3}$, since
$\sum_k\lambda_k^{-4}<\infty$.

We use the realization already proved in
\cref{prop:step-realization}. Let $\Omega_3$ be its weighted
square-summability domain with exponent $p=3$, intersected with
$\Omega$. For every $u\in\Omega_3$, the step distribution has the
absolutely convergent pairings \eqref{eq:absolute-pairing}.
Let $\Omega^{(a)}\subset\Omega_3$ be the set on which it has
finite $H_{-3}$ norm, and set $\sharp_a=0$ off this set.
The measurability argument in \cref{prop:compactness} makes
$\Omega^{(a)}$ measurable and $\sharp_a$ a Borel map into $H_{-4}$.
Every periodic cochain lies in $\Omega_3$ and, by the preceding
finite-volume estimate, in $\Omega^{(a)}$.

For each plaquette $\sigma$, choose a real smooth bump $b^\sigma$
with only its corresponding two-form component nonzero, supported
inside $Q_a(\sigma)$ and of integral $a^{3/2}$. Then
$(b^\sigma)_{[a]}=\mathbf1_\sigma$. The map
\begin{equation}\label{eq:limit-step-inverse}
 \mathcal J_a:H_{-4}\longrightarrow\R^{\mathcal P_2},
 \qquad \mathcal J_a(F)(\sigma)=F(b^\sigma),
\end{equation}
is continuous for the product topology, and
$\mathcal J_a\sharp_a u=u$ on $\Omega^{(a)}$.
At finite volume the identities are
\[
 P_{a,n_j}=(\sharp_a)_*\nu_{a,n_j},\qquad
 \nu_{a,n_j}=(\mathcal J_a)_*P_{a,n_j}.
\]
The mapping theorem therefore gives
$\nu_{a,n_j}\Longrightarrow\nu_a:=(\mathcal J_a)_*\mu_a$.
The closed conditions $du=0$ pass to this product limit.

Fix a finitely supported real $\beta$ on the covering lattice.
For all sufficiently large $n_j$, its support embeds injectively
in the torus. Its folded norm then equals its covering
$\ell^2$ norm. Apply \eqref{eq:recovery-finite-gaussian}, pass
the bounded continuous cylinder
$\min\{R,e^{\langle\beta,u\rangle}\}$ to the limit, and let
$R\to\infty$. This proves \eqref{eq:limit-lattice-G}.

The Gaussian bound implies the finite-coefficient second-moment
bound, so \cref{prop:step-realization} gives $\nu_a(\Omega_3)=1$.
The step distribution on this domain satisfies the common
$L^2$ test estimate. Applying \cref{prop:compactness} with $s=3$
therefore gives $\nu_a(\Omega^{(a)})=1$.
The finite-coefficient Fatou argument in
\cref{cor:qed-continuum} now gives
\[
 \E_{\nu_a}e^{(\sharp_a u)(f)}
 \le e^{\|f_{[a]}\|_{\ell^2}^2/2}
 \le e^{\|f\|_2^2/2}.
\]

It remains to prove the reverse equality in
\eqref{eq:limit-volume-identification}. Choose a countable
Schwartz-dense family $\mathcal D_0$ of real compact smooth
two-forms. For $f\in\mathcal D_0$, the test
\[
 f'=f-\sum_\sigma f_{[a]}(\sigma)b^\sigma
\]
is a finite sum and satisfies $f'_{[a]}=0$. Every finite-volume
step field therefore annihilates $f'$. The hyperplane
$\{F:F(f')=0\}$ is closed in $H_{-4}$ and has full $\mu_a$
measure. Intersect these events for $f\in\mathcal D_0$ with
the full-measure event
$\mathcal J_a(F)\in\Omega^{(a)}$. On that single event,
\[
 F(f)=\sum_\sigma f_{[a]}(\sigma)F(b^\sigma)
      =(\sharp_a\mathcal J_aF)(f)
 \qquad(f\in\mathcal D_0).
\]
Both sides are tempered distributions, so density gives
$F=\sharp_a\mathcal J_aF$ there. Taking laws proves the
pushforward identity and the Gaussian assertion for $\mu_a$.
\end{proof}

\subsection{Finite-torus estimates and joint cutoff sequences}
\label{subsec:joint-cutoff-sequences}
The exact $L^2$ Gaussian bound in
\cref{thm:gaussian-convergence} applies to the recovered
infinite-volume laws. Before the volume limit, the Schwartz
test is periodized. We recall the estimate from
\cite[Sections~6.1--6.2]{QED3Manuscript} that controls this step.

Put $R=an$ and $f^R(y)=\sum_{k\in\Z^3}f(y+Rk)$.
The finite-torus coefficient is
\begin{equation}\label{eq:folded-cell-coefficient}
 b_{a,n}([\sigma])
 =\sum_{k\in\Z^3}f_{[a]}(\sigma+nk)
 =a^{-3/2}\int_{Q_a(\sigma)}f^R_{\ell m}(y)\,\dd y
 \quad(\sigma=(x;\ell m)).
\end{equation}
The sum converges absolutely. It gives
$F_a(f)=\langle b_{a,n},u\rangle$, and cell Cauchy--Schwarz
gives $\|b_{a,n}\|_{\ell^2}\le
\|f^R\|_{L^2(\R^3/R\Z^3)}$. With $Q_R=[-R/2,R/2)^3$, define
\begin{equation}\label{eq:limit-periodic-seminorm}
 p_R(f)=\|f\|_2+
        26R^{-2}\||y|^2f\|_{L^2(\R^3\setminus Q_R)}.
\end{equation}
Then
\begin{equation}\label{eq:limit-cutoff-G}
 \E_{P_{a,n}}e^{F(f)}
 \le e^{\|b_{a,n}\|_{\ell^2}^2/2}
 \le e^{p_R(f)^2/2}
 \qquad(f\text{ real}).
\end{equation}
To see the last estimate, put
$c_k=\|f\|_{L^2(Q_R+Rk)}$. Minkowski gives
$\|f^R\|_{L^2(\R^3/R\Z^3)}\le\sum_kc_k$. For $k\ne0$,
$|y|\ge R|k|_\infty/2$ on $Q_R+Rk$, and weighted
Cauchy--Schwarz gives
\[
 \sum_{k\ne0}c_k
 \le4R^{-2}
   \left(\sum_{k\ne0}|k|_\infty^{-4}\right)^{1/2}
   \||y|^2f\|_{L^2(\R^3\setminus Q_R)}.
\]
The shell count $24j^2+2$ gives
$\sum_{k\ne0}|k|_\infty^{-4}=4\pi^2+\pi^4/45$,
whose square root times four is less than $26$.
For $R\ge1$, one has $p_R(f)\le p_1(f)$, and
$p_R(f)\to\|f\|_2$ as $R\to\infty$. The seminorm $p_1$
is continuous on Schwartz space.

\begin{corollary}\label{cor:prescribed-volume-first}
At each spacing $a_i\downarrow0$, take the chamber torus laws
with parameters fixed during that volume limit. Let $\mu_i$
be any prescribed weak $H_{-4}$ volume limit of their
centered step fields. Suppose $\mu_i\Longrightarrow\mu$ on
$H_{-4}$. Then $\mu$ has the conclusions of
\cref{cor:plaquette-limit}, including convergence of all
finite-family moments and cumulants from $\mu_i$ to $\mu$.
Its tensor Schwinger functions satisfy \cref{prop:tensor-growth}.
There are torus sides $n_i\to\infty$ such that
$a_in_i\to\infty$,
\[
 P_{a_i,n_i}\Longrightarrow\mu,
\]
so $\mu$ is also realized by a joint sequence of finite
cutoff laws.
\end{corollary}

\begin{proof}
Apply \cref{prop:fixed-spacing-recovery} at each $a_i$.
It identifies the prescribed law $\mu_i$ with the step image
of a weak cochain limit. The chamber theorem
\cite[Theorem~A]{ChamberLattice} applies to every such limit
with the parameters fixed: it gives closedness, Gaussian
domination, the lattice symmetries, and reflection positivity
at the nine mirror families. Thus these recovered laws meet
all the hypotheses of \cref{cor:plaquette-limit}; that
corollary and \cref{prop:tensor-growth} give the assertions
about $\mu$ and the sequence $\mu_i$.

Fix a metric for weak convergence of probability laws on
$H_{-4}$. Along the defining volume sequence for $\mu_i$,
choose $n_i$ so large that $a_in_i\ge i$ and the distance
between $P_{a_i,n_i}$ and $\mu_i$ is less than $1/i$.
Then $P_{a_i,n_i}\Longrightarrow\mu$ and
$R_i:=a_in_i\to\infty$. This is the diagonal selection in
\cite[proof of Proposition~6.1]{QED3Manuscript}.
\end{proof}
